% 贯穿图的全部8棵生成树（由脚本枚举5条边中3条边的连通组合得到）。
% 粗蓝线为树边，浅色虚线为舍弃的原边。
\begin{tikzpicture}[>=Stealth,
  v/.style={circle,draw=FigureStroke,fill=FigureBlueFill,line width=.7pt,minimum size=5mm,inner sep=0pt,font=\figurefont\fontsize{7.5}{9}\selectfont},
  t/.style={draw=FigureBlue,line width=1.5pt},
  o/.style={draw=FigureMuted,dashed,line width=.5pt}]
\begin{scope}[xshift=0.00cm,yshift=0.00cm]
  \node[v] (1) at (-0.80,0.00) {$1$};
  \node[v] (2) at (0.00,0.60) {$2$};
  \node[v] (3) at (0.00,-0.60) {$3$};
  \node[v] (4) at (0.80,0.00) {$4$};
  \draw[t] (1)--(2);
  \draw[t] (1)--(3);
  \draw[o] (2)--(3);
  \draw[t] (2)--(4);
  \draw[o] (3)--(4);
  \node at (0,-1.15) {$T_{1}$};
\end{scope}
\begin{scope}[xshift=2.50cm,yshift=0.00cm]
  \node[v] (1) at (-0.80,0.00) {$1$};
  \node[v] (2) at (0.00,0.60) {$2$};
  \node[v] (3) at (0.00,-0.60) {$3$};
  \node[v] (4) at (0.80,0.00) {$4$};
  \draw[t] (1)--(2);
  \draw[t] (1)--(3);
  \draw[o] (2)--(3);
  \draw[o] (2)--(4);
  \draw[t] (3)--(4);
  \node at (0,-1.15) {$T_{2}$};
\end{scope}
\begin{scope}[xshift=5.00cm,yshift=0.00cm]
  \node[v] (1) at (-0.80,0.00) {$1$};
  \node[v] (2) at (0.00,0.60) {$2$};
  \node[v] (3) at (0.00,-0.60) {$3$};
  \node[v] (4) at (0.80,0.00) {$4$};
  \draw[t] (1)--(2);
  \draw[o] (1)--(3);
  \draw[t] (2)--(3);
  \draw[t] (2)--(4);
  \draw[o] (3)--(4);
  \node at (0,-1.15) {$T_{3}$};
\end{scope}
\begin{scope}[xshift=7.50cm,yshift=0.00cm]
  \node[v] (1) at (-0.80,0.00) {$1$};
  \node[v] (2) at (0.00,0.60) {$2$};
  \node[v] (3) at (0.00,-0.60) {$3$};
  \node[v] (4) at (0.80,0.00) {$4$};
  \draw[t] (1)--(2);
  \draw[o] (1)--(3);
  \draw[t] (2)--(3);
  \draw[o] (2)--(4);
  \draw[t] (3)--(4);
  \node at (0,-1.15) {$T_{4}$};
\end{scope}
\begin{scope}[xshift=0.00cm,yshift=-2.30cm]
  \node[v] (1) at (-0.80,0.00) {$1$};
  \node[v] (2) at (0.00,0.60) {$2$};
  \node[v] (3) at (0.00,-0.60) {$3$};
  \node[v] (4) at (0.80,0.00) {$4$};
  \draw[t] (1)--(2);
  \draw[o] (1)--(3);
  \draw[o] (2)--(3);
  \draw[t] (2)--(4);
  \draw[t] (3)--(4);
  \node at (0,-1.15) {$T_{5}$};
\end{scope}
\begin{scope}[xshift=2.50cm,yshift=-2.30cm]
  \node[v] (1) at (-0.80,0.00) {$1$};
  \node[v] (2) at (0.00,0.60) {$2$};
  \node[v] (3) at (0.00,-0.60) {$3$};
  \node[v] (4) at (0.80,0.00) {$4$};
  \draw[o] (1)--(2);
  \draw[t] (1)--(3);
  \draw[t] (2)--(3);
  \draw[t] (2)--(4);
  \draw[o] (3)--(4);
  \node at (0,-1.15) {$T_{6}$};
\end{scope}
\begin{scope}[xshift=5.00cm,yshift=-2.30cm]
  \node[v] (1) at (-0.80,0.00) {$1$};
  \node[v] (2) at (0.00,0.60) {$2$};
  \node[v] (3) at (0.00,-0.60) {$3$};
  \node[v] (4) at (0.80,0.00) {$4$};
  \draw[o] (1)--(2);
  \draw[t] (1)--(3);
  \draw[t] (2)--(3);
  \draw[o] (2)--(4);
  \draw[t] (3)--(4);
  \node at (0,-1.15) {$T_{7}$};
\end{scope}
\begin{scope}[xshift=7.50cm,yshift=-2.30cm]
  \node[v] (1) at (-0.80,0.00) {$1$};
  \node[v] (2) at (0.00,0.60) {$2$};
  \node[v] (3) at (0.00,-0.60) {$3$};
  \node[v] (4) at (0.80,0.00) {$4$};
  \draw[o] (1)--(2);
  \draw[t] (1)--(3);
  \draw[o] (2)--(3);
  \draw[t] (2)--(4);
  \draw[t] (3)--(4);
  \node at (0,-1.15) {$T_{8}$};
\end{scope}
\end{tikzpicture}
